\documentclass[11pt]{article}
\usepackage[margin=1in]{geometry}
\usepackage{amsmath,amssymb,amsthm,mathtools}
\usepackage{graphicx}
\usepackage[hidelinks]{hyperref}
\usepackage{microtype}

\newtheorem{theorem}{Theorem}
\newtheorem{lemma}{Lemma}
\newtheorem{remark}{Remark}
\newcommand{\M}{\mathcal M}
\newcommand{\PiD}{\Pi_{\delta}}
\newcommand{\dd}{\,d}

\title{The relative $p$-Wasserstein distance is a metric\\
and its two coupling formulas agree}
\author{Candidate solution packet for arXiv:2411.05678}
\date{21 July 2026}

\begin{document}
\maketitle

\begin{center}
\fbox{\parbox{0.92\textwidth}{
\textbf{Status.} Candidate full solution, likely valid, pending human review.
For every $1\leq p<\infty$, the packet proves $\widehat W_p=W_p$ and proves
identity of indiscernibles for $W_p$.  Thus it answers both questions posed
for $p>1$ in Section 6.3 of the source paper.
}}
\end{center}

\section{Source questions}

The source is P.~Bubenik and A.~Elchesen, \emph{Relative Optimal Transport},
arXiv:2411.05678v3 (6 April 2026) \cite{BE}.  In Section 6.3, on source PDF
pages 31 and 32, the authors ask:
\begin{quote}
``It is an open question whether $\widehat W_p=W_p$ for $p>1$.''

``It is an open question whether $W_p(\mu,\nu)=0$ implies that $\mu=\nu$
when $p>1$.''
\end{quote}

\begin{figure}[ht]
  \centering
  \includegraphics[width=0.98\textwidth]{figures/open_problem_crop.png}
  \caption{The two questions on source PDF pages 31 and 32.}
\end{figure}

\section{Definitions and notation}

Let $(X,d,A)$ be a metric pair as in the source: $(X,d)$ is complete and
separable, and $A\subset X$ is nonempty and closed.  Put
\[
 d_A(x):=d(x,A),\qquad
 d_p(x,y):=\min\!\left\{d(x,y),
       \bigl(d_A(x)^p+d_A(y)^p\bigr)^{1/p}\right\}.
\]
A relative Borel measure is a Borel measure modulo measures supported on
$A$.  Let $\M_p^+(X,A)$ be the source's space of tight relative measures
$\mu$ with
\[
 \mu(d_A^p)=\int_{X\setminus A}d_A(x)^p\dd\mu(x)<\infty.
\]
For $\delta>0$, write $\mu_\delta$ for the restriction to
$\{d_A>\delta\}$ and $\mu^\delta$ for the restriction to
$\{0<d_A\leq\delta\}$.  The source proves
\begin{equation}
 \mu^\delta(d_A^p)\longrightarrow0 \label{eq:tails}
 \quad\text{as }\delta\downarrow0.
\end{equation}

Let $\Pi(\mu,\nu)$ be the set of relative couplings.  Fix $a\in A$.
The structured class $\PiD(\mu,\nu)$ consists of the trivial extensions
\[
 \rho=\sigma+(i_1^a)_*\mu^\delta+(i_2^a)_*\nu^\delta,
 \qquad \sigma\in\Pi(\mu_\delta,\nu_\delta),
\]
where $i_1^a(x)=(x,a)$ and $i_2^a(y)=(a,y)$.  The two quantities in the
first open question are
\begin{align*}
 \widehat W_p(\mu,\nu)^p
   &:=\inf_{\pi\in\Pi(\mu,\nu)}\int d_p(x,y)^p\dd\pi(x,y),\\
 W_p(\mu,\nu)^p
   &:=\inf_{\delta>0}\ \inf_{\rho\in\PiD(\mu,\nu)}
       \int d_p(x,y)^p\dd\rho(x,y).
\end{align*}
Since every structured coupling is a coupling, the source observes that
$\widehat W_p\leq W_p$.

\section{Proof intuition}

The first question is resolved by approximating an \emph{arbitrary} coupling,
not an optimizer.  Collapse only the $\delta$-neighborhood of the reservoir
$A$ onto one point of $A$, and then restore the two collapsed marginal tails
by the trivial reservoir couplings.  The resulting coupling is structured.
The collapsed integrand converges pointwise to the original integrand and is
dominated by
$d_A(x)^p+d_A(y)^p$, whose coupling integral is exactly the sum of the two
finite marginal $p$-moments.  Dominated convergence, together with
\eqref{eq:tails},
therefore preserves the cost in the limit.  This gives the missing inequality
$W_p\leq\widehat W_p$ without needing a quantitative $p>1$ retraction
estimate.

For definiteness, take a compact set $K$ away from $A$.  A particle starting
in $K$ and ending outside a small ordinary $d$-neighborhood $K_r$ must pay at
least $r$ in the relative metric $d_p$: both the direct branch and the
reservoir branch in the definition of $d_p$ exceed $r$.  A coupling of
arbitrarily small cost therefore gives $\mu(K)\leq\nu(K_r)$.  The set $K_r$
has finite $\nu$-mass because it remains a positive distance from $A$.
Continuity from above as $r\downarrow0$, followed by symmetry and tightness,
then forces $\mu=\nu$.  No properness or compactness of $K_r$ is used.

\section{The full result}

\begin{theorem}\label{thm:main}
Let $(X,d,A)$ be a complete separable metric pair with $A\neq\varnothing$,
and let $1\leq p<\infty$.  For all
$\mu,\nu\in\M_p^+(X,A)$,
\[
             W_p(\mu,\nu)=\widehat W_p(\mu,\nu).
\]
Moreover,
\[
             W_p(\mu,\nu)=0\quad\Longrightarrow\quad\mu=\nu.
\]
Consequently the pseudometric $W_p$ constructed in the source is a metric on
$\M_p^+(X,A)$.  In particular, both open questions stated there for $p>1$
have affirmative answers.
\end{theorem}

\begin{proof}
We first prove equality of the two costs.  Fix a coupling
$\pi\in\Pi(\mu,\nu)$.  For $\delta>0$, define the Borel map
\[
 q_\delta(x):=
 \begin{cases}
   a,&0<d_A(x)\leq\delta,\\
   x,&\text{otherwise},
 \end{cases}
 \qquad R_\delta(x,y):=(q_\delta(x),q_\delta(y)).
\]
This is the retraction used in Section 6.2 of the source.  Its coupling lemma
gives
\[
       \sigma_\delta:=(R_\delta)_*\pi
       \in\Pi(\mu_\delta,\nu_\delta).
\]
Trivially extend it to
\[
 \rho_\delta:=\sigma_\delta
       +(i_1^a)_*\mu^\delta+(i_2^a)_*\nu^\delta
       \in\PiD(\mu,\nu).
\]
Because $d_p(x,a)=d_A(x)$, its cost is exactly
\begin{align}
 \int d_p^p\dd\rho_\delta
 &=\int_{X^2}d_p(q_\delta(x),q_\delta(y))^p\dd\pi(x,y) \notag\\
 &\quad+\mu^\delta(d_A^p)+\nu^\delta(d_A^p). \label{eq:cost}
\end{align}
For every $x\in X$, $q_\delta(x)\to x$ as $\delta\downarrow0$.
Continuity of $d_p$ therefore gives pointwise convergence of the integrand
in \eqref{eq:cost}.  Furthermore,
\begin{align}
 d_p(q_\delta(x),q_\delta(y))^p
 &\leq d_A(q_\delta(x))^p+d_A(q_\delta(y))^p \notag\\
 &\leq d_A(x)^p+d_A(y)^p. \label{eq:dominator}
\end{align}
The last function is $\pi$-integrable, since the coupling marginals give
\begin{equation}
 \int_{X^2}\bigl(d_A(x)^p+d_A(y)^p\bigr)\dd\pi(x,y)
   =\mu(d_A^p)+\nu(d_A^p)<\infty. \label{eq:moments}
\end{equation}
Dominated convergence and \eqref{eq:tails} applied to \eqref{eq:cost} now yield
\[
       \lim_{\delta\downarrow0}\int d_p^p\dd\rho_\delta
       =\int d_p^p\dd\pi.
\]
Since every $\rho_\delta$ is admissible in the definition of $W_p$,
\[
       W_p(\mu,\nu)^p\leq\int d_p^p\dd\pi.
\]
Taking the infimum over arbitrary $\pi\in\Pi(\mu,\nu)$ gives
$W_p\leq\widehat W_p$.  Together with the reverse inequality already
noted in the source, this proves $W_p=\widehat W_p$.

We next prove identity of indiscernibles directly.  Assume
$W_p(\mu,\nu)=0$.  Let $K\subset X\setminus A$ be compact and put
\[
        \alpha:=d(K,A)>0.
\]
Fix $0<r<\alpha$ and let
\[
        K_r:=\{y\in X:d(y,K)\leq r\}.
\]
The function $d_A$ is $1$-Lipschitz, so
$d_A(y)\geq\alpha-r$ on $K_r$.  Hence
\begin{equation}
       \nu(K_r)\leq(\alpha-r)^{-p}\nu(d_A^p)<\infty. \label{eq:finite-neighborhood}
\end{equation}
For $x\in K$ and $y\notin K_r$, one has $d(x,y)>r$, while
\[
 \bigl(d_A(x)^p+d_A(y)^p\bigr)^{1/p}
       \geq d_A(x)\geq\alpha>r.
\]
Thus both branches defining $d_p(x,y)$ exceed $r$, and consequently
\begin{equation}
       d_p(x,y)>r \label{eq:separation}
       \qquad(x\in K,\ y\notin K_r).
\end{equation}

Given $\eta>0$, the definition of $W_p=0$ supplies some $\delta>0$ and
$\rho\in\PiD(\mu,\nu)$ such that
$\int d_p^p\dd\rho<\eta$.  Marginality, \eqref{eq:separation}, and Markov's inequality give
\begin{align*}
 \mu(K)=\rho(K\times X)
 &\leq \rho(K\times K_r)+\rho(K\times(X\setminus K_r))\\
 &\leq \nu(K_r)+r^{-p}\int d_p^p\dd\rho\\
 &<\nu(K_r)+r^{-p}\eta.
\end{align*}
Letting $\eta\downarrow0$ gives $\mu(K)\leq\nu(K_r)$.
Choose any decreasing sequence $r_n\downarrow0$ with $r_1<\alpha$.
Then $K_{r_n}\downarrow K$, and \eqref{eq:finite-neighborhood} makes the first set finite.  Continuity
from above therefore yields
\[
                  \mu(K)\leq\nu(K).
\]
Interchanging $\mu$ and $\nu$ gives equality on every compact
$K\subset X\setminus A$.  Both measures are tight by the definition of
$\M_p^+(X,A)$.  The source's compact-set uniqueness lemma for tight relative
measures now implies $\mu=\nu$.

The source already proves reflexivity, symmetry, and the triangle inequality
for $W_p$ on $\M_p^+(X,A)$.  Adding identity of indiscernibles completes the
metric assertion.
\end{proof}

\section{Verification report}

The proof is analytic and no numerical evidence is used.  The following
points were checked against the definitions and lemmas in the current source:
\begin{itemize}
  \item The retraction is Borel and its pushforward has marginals
  $\mu_\delta,\nu_\delta$ (source Lemma 6.21).
  \item Formula \eqref{eq:cost} is exact: each tail is coupled to a point of $A$, where
  $d_p(x,a)^p=d_A(x)^p$.
  \item The dominator in \eqref{eq:dominator} is integrable even when the coupling has infinite
  total mass, because \eqref{eq:moments} depends only on the two finite $p$-moments.
  \item In the definiteness proof, $K_r$ need not be compact.  Its finite mass
  follows from \eqref{eq:finite-neighborhood}, which is sufficient for continuity from above.
  \item All marginal identities are evaluated on sets disjoint from $A$, so
  they are independent of the representative chosen for a relative measure.
  \item Agreement on compact subsets of $X\setminus A$ implies equality by
  the source's Lemma 4.1 and the tightness assumption.
\end{itemize}

The main human-review focus should be the passage from a relative coupling to
the exact integral identity \eqref{eq:cost}, including representative-independence, and
the compact-set limiting argument following \eqref{eq:separation}.  Current verdict:
\textbf{full solution likely valid}.

\section{Scope and novelty check}

The theorem uses exactly the source hypotheses: a complete separable metric
pair, a nonempty closed reservoir, and tight relative measures with finite
$p$-moment.  It does not require an optimal coupling, finite total mass,
properness of $X$, or compactness of metric neighborhoods.  Although the
source questions are stated for $p>1$, the argument works uniformly for every
$1\leq p<\infty$; for $p=1$ it recovers conclusions already proved in the
source.

Before promotion, the local registry, solution, attempt, and proof-gap indexes
were searched for arXiv:2411.05678 and the core phrases ``relative optimal
transport'', ``relative Wasserstein'', ``$\widehat W_p=W_p$'', and ``identity
of indiscernibles''; no prior result was found.  A bounded web/arXiv search on
21 July 2026 used the title, authors, exact open-question phrases, and nearby
relative Wasserstein terminology.  It found the source and adjacent optimal
transport work but no later settlement.  The current source version, v3 of
6 April 2026, still labels both questions open.  This supports, but does not
establish, novelty; the search was not an exhaustive review of all published
literature.

\begin{thebibliography}{9}
\bibitem{BE}
P.~Bubenik and A.~Elchesen,
\emph{Relative Optimal Transport}, arXiv:2411.05678v3, 2026.
\end{thebibliography}

\end{document}
